% Pursuing Stacks (cotes 134-2 to 134-8): archive cote 134-4, site page 11 (PDF page 12), typescript page 99.
% typescript. Draft compared: ps3.tex:537-605.
% AI-transcribed by Claude Opus 5.5 (claude-opus-5-5, adaptive thinking, effort high), Message Batches API, batch ps1 (msgbatch_01Rv5hCaNTtUt4JHx1yqoHPL), 2026-11-12; not yet checked by eye.
\scanpage{134-4/11}
\mspage{\hand{99}}
Let X,Y be objects of M,  and \xout{f,g}
\[ f,g:\ X\rightrightarrows Y \]
two maps in M from X to Y. One key notion constantly used lately (but so
far only when f is an identity map,  and g a"constant"one  -  which is
equally the case needed for \inserted[typed]{defining} \xout{dealing with} contractibility) is the notion
of an $\underline{I}$-\uline{homotopy} \xout{between} \uline{from} f \uline{to} g,  namely a map $F\mathrm{x}I\xrightarrow{h}G$  making com-%
mutative the diagram
\diagram{\begin{tikzcd}
 & F\mathrm{x}I \ar[dd,"\hand{h}"] & \\
F \ar[ur,"\hand{\mathrm{id}_F\times\delta_0}"] \ar[dr,swap,"\hand{f}"] & & \over{F}{\unsure{F}} \ar[ul,swap,"\hand{\mathrm{id}_F\times\delta_1}"] \ar[dl,"\hand{g}"] \\
 & G &
\end{tikzcd}} .

Let's first restate the "homotopy lemma" of page 60 in a slightly more
complete form:

\uline{Homotopy lemma reformulated}:  Assume f and g are $\underline{I}$-homotopic.  Then:
a) $\hand{\gamma}(f)=\hand{\gamma}(g)$,  where
\[ \hand{\gamma}:\ M\longrightarrow W^{-1}M = H_{M} \]
is the canonical functor.

b) If  f  is a weak equivalence,  so is g  (and conversely of course,
by symmetry\xout{ix} of the roles of $\hand{\delta}_0$ and $\hand{\delta}_1$).

c) Assume f is \inserted[typed]{an isomorphism,  and g constant  -  then  F} \xout{\unsure{W}-aspheric, then so is g - and conversely} $\hand{\longrightarrow}$ e  and
$\overset{\hand{\unsure{\hat{\ast}}}}{G}\rightarrow e^{\hand{\ast}}$  are W-aspheric.

It is important to notice that the \xout{notion} relation `` f is $\underline{I}$-homotopic
to g'' in Hom(X,Y) is not necessarily symmetric nor transitive,  and that
it depends on the choice of the homotopy interval ($\underline{I}$,$\hand{\delta}_0$,$\hand{\delta}_1$).  Thus the
symmetric relation from  $\underline{I}$-homotopy is $\hand{\check{\phantom{I}}}\!\!\underline{I}$-homotopy,  where  $\hand{\check{\phantom{I}}}\!\!\underline{I}$ is the
homotopy interval "opposite" to $\underline{I}$  (namely with $\hand{\delta}_0$,$\hand{\delta}_1$ reversed).  The
example we are immediately interested in is  M = (Cat),  with $W=W_{(Cat)}$ the
usual notion of weak equivalence.  A homotopy interval is just an aspheric
small category I,  endowed with two \uline{distinct} objects  $e_0$ , $e_1$.(The condi-%
tion $e_0\hand{\neq}e_1$  just expresses the condition\inserted{b)} $e_0\wedge e_1 = \emptyset$  on homotopy intervals
-  if it were not satisfied,  $\underline{I}$-homotopy would just mean \inserted[typed]{equality} \xout{identity} of f and
g ...).  For the usual choice  $I =\hand{\Delta}_1$ =  ($e_0\hand{\rightarrow} e_1$),  an\deleted{y} $\underline{I}$-homotopy from
f to g is just a  morphism  between functors $f\rightarrow g$  -  the I-homotopy rela-%
tion between f and g is the existence of such a morphism,  it is atransiti
ve\hand{,}and \inserted[typed]{generally} non-symmetric relation.  If we take I to be a category with \over{\unsure{g}}{j}ust
two objects $e_0$ and $e_1$  ,  equivalent to the final category ,  an I-homoto-%
py between f and g is  just an isomorphism from f to g  -  the I-homotopy
relation now is both transitive and symmetric,  and it is a lot more res-%
trictive than the previous one.  If we take I to be the barycentric subdi-%

\margin{11}
